% ---------------------------------------------
% Copyright Ben Paul Wise. All rights reserved.
% ---------------------------------------------
% SBC versions 1 and 2 and the FPMQ2 and FPMQ3 password files.
% Process with pdfLaTeX, then BibTeX, then pdfLaTeX twice. The references come
% from CPVI/shared-tex/fdp.bib.

\input{../../../CPVI/shared-tex/prolog-article.tex}

\usepackage{array} % ragged-right table columns
\usepackage{placeins} % \FloatBarrier keeps the algorithms in their section
\setlength{\headheight}{14.5pt} % room for the 12 pt running head

% Start pages that hold only floats (the algorithms) at the top, not centred,
% and let tall algorithms share a page with text or with each other.
\makeatletter
\setlength{\@fptop}{0pt}
\makeatother
\renewcommand{\topfraction}{0.9}
\renewcommand{\bottomfraction}{0.8}
\renewcommand{\textfraction}{0.1}
\renewcommand{\floatpagefraction}{0.8}

\lhead{B. P. Wise}
\rhead{SBC version 2 and FPMQ3}
\cfoot{\thepage\ of \pageref{LastPage}}

\newcommand{\code}[1]{\texttt{#1}}
\newcommand{\xor}{\oplus}
\newcommand{\rotl}{\lll}

\title{Sliding Block Cipher Version 2\\ and the FPMQ3 Password File:\\
Design, Repairs, and Open Questions}
\author{Ben P. Wise}
\date{21 October 2026}

\begin{document}

\maketitle

\begin{abstract}
The Sliding Block Cipher (SBC) is a block cipher of the author's design, first
built in 2002 \cite{wise2006}. It encrypts the password file of \code{fpwdman}, the author's
password manager. A review in October 2026 found no practical way to break SBC,
but it found one structural weakness in the cipher and four weaknesses in the
file format and the program around it. This article explains in plain language how the cipher
and the file format work and why each part is there. It then describes the
weaknesses and how they were repaired in SBC version 1 and in a new file format
called FPMQ2, which \code{fpwdman} 4.0 writes. Two further changes, a rotation
whose amount depends on the data and a permutation that depends on the key,
make SBC version 2 and the file format FPMQ3, which \code{fpwdman} 4.1 writes.
The article gives pseudo-code for the cipher and the file format, and it lists
the attacks and analyses that remain to be done.
\end{abstract}


\clearpage
\tableofcontents


\clearpage
\section{Introduction}

A password manager keeps many secrets under one secret. The user remembers a
single master passphrase; the program keeps every other password in a file
encrypted under a key made from that passphrase. Anyone who copies the file can
try to open it at leisure, on as many computers as they can afford. The design
must therefore make two things hard: recovering the contents without the
passphrase, and finding the passphrase by guessing.

\code{fpwdman} meets the first need with the Sliding Block Cipher (SBC), and the
second with a standard key-derivation function. Section \ref{sec:how} explains
both in plain language. Section \ref{sec:fixed} describes the weaknesses found
in the review and how they were repaired. Section \ref{sec:pseudo} gives
pseudo-code for the cipher and for the file format. Section \ref{sec:open}
lists what remains to be studied. Appendix \ref{sec:measure} records the
measurements on which the article relies.

Four file formats appear throughout:
\begin{itemize}[nosep]
\item the \emph{legacy} format, written by the original FLTK program, which is
      read but no longer written;
\item \emph{FPMQ1}, written by the Qt program before version 4.0, which is also
      read but no longer written;
\item \emph{FPMQ2}, written by version 4.0, which is also read but no longer
      written;
\item \emph{FPMQ3}, written by version 4.1 and later.
\end{itemize}
Opening a file in an older format and saving it converts it to FPMQ3. The
cipher likewise has three versions: version 0 for the legacy format and FPMQ1,
version 1 for FPMQ2, and version 2 for FPMQ3.

\section{How it works, and why}
\label{sec:how}

\subsection{The cipher}

SBC treats a block of data as a ring of $K$ words of 32 bits each. The
\code{fpwdman} settings use $K = 32$, so a block holds 1024 bits (128 bytes).
The ring has no first or last word: the word before word 0 is word $K-1$.
A classic Feistel cipher (such as the US Data Encryption Standard, Blowfish, and many more)
uses the same ring structure with $K=2$. As there are only two elements in such a  ring,
they are conventionally referred to as the left and right blocks.

Encryption proceeds one \emph{step} at a time. A step changes exactly one word,
using the word just before it in the ring and one 32-bit \emph{subkey}. Step 0
changes word 0, step 1 changes word 1, and so on around the ring. One pass
around the ring is one \emph{round}. \code{fpwdman} uses $R = 24$ rounds, that
is $24 \times 32 = 768$ steps.

The change made at each step is
\begin{eqnarray}
q(x) &=& (x+1)(2x+1) \bmod 2^{32}, \label{eq:step_A} \\
x_p &\leftarrow& x_p + \bigl( q(x_{p-1} \xor s_i) \rotl 7 \bigr) \bmod 2^{32},
\label{eq:step_B}
\end{eqnarray}
where $x_p$ is the word being changed, $x_{p-1}$ the word before it, $s_i$ the
subkey for step $i$, $\xor$ is bitwise exclusive-or, and $\rotl 7$ rotates the
32 bits left by seven places. The function $q$ (called \code{qTransform} in the
code) is a quadratic that maps each 32-bit number to a different 32-bit number.
It was adapted from the cipher RC6, and the added constant 1 removes the fixed
points that RC6's version has.

The author adapted the quadratic transformation of RC6 to 
derive the general form of $q$ as follows.

\begin{equation}
q(x) = (x+a)(nx+c) \bmod 2^{m},
\end{equation}

To ensure that  $q(x) = q(y)$ if and only if $x = y$ , we require
that 
$n$ is even positive, 
and
$c$ be odd positive.  We prove the bilateral implication first by assuming $x \neq y$
and showing that $q(x) \neq q(y)$ follows.
\begin{eqnarray}
q(x) - q(y) &=& \left( nx^2 +(an+ c)x + ac \right) - \left( ny^2 + (an+c)y  + ac \right) \label{eq:qinj_A} \\
 &=& \left( n (x+y) + (an+c) \right)  (x-y) \label{eq:qinj_B}
\end{eqnarray}

Here $x$ and $y$ are words, numbers from $0$ to $2^m - 1$, and the right-hand
side is computed in ordinary integers. Without loss of generality, assume that $x > y$. Write $u$ for 
the first factor of the right-hand side:
\begin{eqnarray}
u &=& n(x+y) + an + c. \label{eq:qinj_C}
\end{eqnarray}
The number $u$ is odd, because $n(x+y)$ and $an$ are even (as $n$ is even) and
$c$ is odd.

Suppose, to the contrary, that $q(x) = q(y)$, that is, that $q(x) - q(y)$ is a
multiple of $2^m$. Then $2^m$ divides the product $u\,(x-y)$. Since $u$ is odd,
it contains no factor of 2, so all $m$ of the factors of 2 must lie in $x - y$:
\begin{eqnarray}
2^m \mid u\,(x-y) \ \text{ and } \ u \text{ odd} &\Longrightarrow& 2^m \mid (x-y).
\label{eq:qinj_D}
\end{eqnarray}
But $x$ and $y$ are different numbers from $0$ to $2^m - 1$, so
$0 < x - y < 2^m$, and no multiple of $2^m$ lies strictly between $0$ and
$2^m$. The supposition is therefore false, and $q(x) \neq q(y)$.

The converse, that $x = y$ gives $q(x) = q(y)$, is immediate. Hence
$q(x) = q(y)$ if and only if $x = y$: $q$ maps the $2^m$ possible words to
$2^m$ different words, so it merely reorders them, and it can be undone.

To ensure that  $q$ has no fixed points , we require
that  $a$ is odd positive. The proof is particularly simple.
Expanding the product and subtracting $x$ gives
\begin{eqnarray}
q(x) - x &=& \left( n x^2 + (an+c)\,x + ac \right) - x \label{eq:qfix_A} \\
 &=& n x^2 + (an + c - 1)\,x + ac. \label{eq:qfix_B}
\end{eqnarray}
The first term is even, because $n$ is even. The second is even, because $an$
is even and $c - 1$ is even, $c$ being odd. The last term, $ac$, is odd,
because $a$ and $c$ are both odd. The right-hand side is therefore odd for
every $x$. But $q(x) = x$ would require $q(x) - x$ to be a multiple of $2^m$,
and so even. Hence there can be no fixed point.

The choice of $n$ matters in one further way: it should be twice an odd
number. Expanding the product gives
\begin{eqnarray}
q(x) &=& n x^2 + (an+c)\,x + ac \bmod 2^{m}. \label{eq:qparts_A}
\end{eqnarray}
Let $2^k$ be the largest power of two that divides $n$. The squared term
$n x^2$ is then zero modulo $2^k$, so the lowest $k$ bits of $q(x)$ come from the
other two terms alone:
\begin{eqnarray}
q(x) &\equiv& (an+c)\,x + ac \pmod{2^k}. \label{eq:qparts_B}
\end{eqnarray}
In those $k$ bits, $q$ only multiplies its input by an odd number and adds a
constant, which is close to linear; the nonlinear squared term begins only at
bit $k$. It begins as low as it can, at bit 1, when $k = 1$, that is, when $n$
is twice an odd number. The present choice $n = 2$ meets this. A multiple of
four does not: with $(a,n,c) = (3,8,17)$, for example, $k = 3$, and the lowest
three bits of $q(x)$ are simply those of $x + 3$. Whether the odd factor is
prime does not matter; only its oddness does. (Four is twice a prime, and fails.)
A larger $n = 2u$, with $u$ odd and with many of its bits set, would still
satisfy the rule and might spread each input bit into more of the higher output
bits in a single application; whether it actually does would have to be
measured.

The final set of conditions are that 
$a$ is odd,
$n$ is twice an odd number, 
$c$ be odd,
and all three are positive.

In version 2 each step ends with a second rotation, of the word being changed,
by an amount taken from the data. With $t$ the output of $q$ and $r$ its five
highest bits, a number from 0 to 31, the step is
\begin{eqnarray}
t &=& q(x_{p-1} \xor s_i), \label{eq:stepv2_A} \\
r &=& \lfloor t / 2^{27} \rfloor, \label{eq:stepv2_B} \\
x_p &\leftarrow& \bigl( x_p + (t \rotl 7) \bmod 2^{32} \bigr) \rotl r.
\label{eq:stepv2_C}
\end{eqnarray}
Versions 0 and 1 stop at step (\ref{eq:step_B}), which is (\ref{eq:stepv2_C})
with $r = 0$.

Every 64 steps (two rounds) the words of the ring are rearranged by a
permutation called \code{perm1}. It does not change any word; it only moves
words to new places. The new order depends on the data: \code{perm1} combines
all the words into one 32-bit number, by applying $q$ to each and taking the
exclusive-or of the results, and uses that number to choose the rearrangement.
It does so in eight moves. At each it adds a different part of the number to
every word's position, modulo $K$, and then folds the positions over: a word in
the first half moves to twice its position, and a word in the second half moves,
in reverse order, to an odd position. Version 2 makes two changes. It combines the number, by exclusive-or,
with one more secret subkey for each permutation, so that the rearrangement can
no longer be computed without the key. And it rotates the number by five places
between the moves, instead of by seven, so that the eight moves use all 32 of
its bits; with seven, two of the bits are never used.

A file is longer than one block, so the blocks are chained: each block is
combined, by exclusive-or, with the encrypted block before it, and then
encrypted. This is the standard cipher-block-chaining (CBC) mode. Decryption
runs every step backward, from the last step to the first, subtracting where
encryption added.

\subsection{Why each part of the cipher is there}

\begin{description}[style=nextline]
\item[The step can always be undone.]
Step (\ref{eq:step_B}) changes $x_p$ by an amount that depends only on
$x_{p-1}$ and $s_i$, and $x_{p-1}$ is not changed by that step. Whoever holds
the key can therefore compute the same amount again and subtract it. This is
the same idea as the Feistel ciphers, generalized from two words to a ring of
any size, which is what lets SBC handle blocks of any length.

\item[The step mixes bits nonlinearly.]
Encryption that used only exclusive-or and addition could be undone by solving
linear equations. The product in $q$ and the mixing of exclusive-or with
addition prevent that.

\item[The rotation moves the high bits down.]
In $q$, each output bit depends only on the input bits at the same place and
below. Without the rotation, a change in a high bit could never reach the low
bits. Rotating by seven places brings high bits to the bottom, so over many
steps every bit influences every other.

\item[The second rotation depends on the data (version 2).]
The amount $r$ comes from the five highest bits of $q$, which depend on every
input bit. Where each bit of the changed word lands therefore depends on the
whole of $x_{p-1}$ and the subkey, so within a single step the low bits of the
result depend on the high bits of the input. The fixed rotation by seven is
kept, so that the high bits of $q$ still reach the bottom in the one step in 32
where $r$ is zero. The step can still be undone, because $r$, like $t$, depends
only on $x_{p-1}$ and $s_i$. The rotation is computed by a single processor
instruction whose time does not depend on the amount, so the time a step takes
reveals nothing about $r$.

\item[The ring spreads a change to every word.]
A change to word $p$ alters the next step's change to word $p+1$, which alters
word $p+2$, and so on. Measured on random data, one round changes about a
quarter of the 1024 bits of a block, and two rounds change half of them, which
is what a random result would do (Appendix \ref{sec:measure}).

\item[The permutation moves words apart.]
The step relates only neighbouring words. Moving words to new places between
rounds makes words that were far apart into neighbours and vice versa.

\item[The permutation can be undone, and in version 2 it is secret.]
The number that chooses the rearrangement does not depend on the order of the
words, so decryption can compute it again from the rearranged block. In version
2 it also includes a subkey, so the rearrangement cannot be computed without the
key. Because the number depends on every word, the two blocks of a pair that
differ almost always receive different rearrangements, which breaks the pattern
of differences that a differential attack follows.

\item[Many rounds give a margin.]
Most attacks on block ciphers grow harder with every round. Twenty-four rounds
are far more than are needed to mix the bits, which leaves room against attacks
not yet known.

\item[Every step has its own subkey (versions 1 and 2).]
Section \ref{sec:slide} explains why this matters.
\end{description}

\subsection{The password file}

Saving the password file in FPMQ3 takes the following stages. (FPMQ2 differs
only in using SBC version 1 and its own labels in stage 3.)
\begin{enumerate}[nosep]
\item Choose 16 random bytes, the \emph{salt}.
\item Run PBKDF2 \cite{rfc8018} on the passphrase and salt, with 600,000
      repetitions of HMAC-SHA-256 \cite{rfc2104,fips180}, to obtain a 32-byte
      \emph{master key}.
\item Expand the master key with HKDF \cite{rfc5869} into two independent keys:
      64 bytes for the cipher and 32 bytes for the authentication code.
\item Place 16 more random bytes, the \emph{nonce}, before the list of entries,
      followed by its length; pad the whole to a multiple of 128 bytes.
\item Encrypt it with SBC version 2 in CBC mode.
\item Compute an authentication code (HMAC-SHA-256) over the file's header and
      the encrypted text, and write the header, the code, and the encrypted
      text.
\end{enumerate}
Opening the file repeats stages 2 and 3, recomputes the authentication code,
and decrypts only if it matches.

\subsection{Why each part of the file is there}

\begin{description}[style=nextline]
\item[The 600,000 repetitions make guessing slow.]
The user pays the cost once when opening the file, about a second. An attacker
pays it for every passphrase tried. This cost is the main protection for a
stolen file, and it must not be reduced to make the program feel faster.

\item[The salt prevents precomputation.]
Because every file has its own salt, an attacker cannot compute a table of
common passphrases and their keys once and use it against every file. Each
guess is good for one file only. The salt need not be secret, and it cannot be:
the program needs it to make the key, and the key is needed to decrypt anything.
It is therefore written in the clear, as in every standard format.

\item[The nonce makes every encryption different.]
The random bytes at the front of the plaintext change the whole encrypted file,
because CBC carries the change from the first block into every later one. Since
the salt is also new at every save, the key changes too; the nonce is a second
safeguard.

\item[The authentication code detects a wrong passphrase and any tampering.]
The code is checked before anything is decrypted. A wrong passphrase gives a
wrong key, and therefore a code that does not match. A file changed by anyone
without the passphrase also fails the check. Checking the code before
decrypting (``encrypt-then-MAC'') is the order that is known to be safe
\cite{bellare2000}.

\item[The cipher and authentication keys are separate.]
Using one key for two purposes invites unforeseen interactions between them.
\end{description}

\section{Weaknesses found and repaired}
\label{sec:fixed}

Table \ref{tab:fixes} summarizes the weaknesses. The subsections explain each.

\begin{table}[ht]
\centering
\begin{tabular}{>{\raggedright\arraybackslash}p{0.24\textwidth}
                >{\raggedright\arraybackslash}p{0.30\textwidth}
                >{\raggedright\arraybackslash}p{0.36\textwidth}}
\hline
\textbf{Weakness} & \textbf{Effect} & \textbf{Repair} \\
\hline
Every round of SBC the same &
Slide attacks, whose cost does not depend on the number of rounds &
SBC version 1: a different subkey for every step, made with HKDF \\
Iteration count read from the file without a check &
A crafted file could freeze the program for about 40 minutes &
Only 600,000 is accepted, checked before any key is made \\
Three PBKDF2 runs per open &
The user paid three times what an attacker pays &
FPMQ2: one PBKDF2 run, expanded with HKDF \\
Header outside the authentication code &
Header could be changed without detection &
FPMQ2: the code covers the header \\
A backup copy kept at every save &
The previous file survived in its old format and under its old passphrase &
Version 4.0 keeps no backup \\
A fixed rotation of seven places &
Low output bits depend on high input bits only after several steps; the step
keeps the regular algebra of a quadratic &
SBC version 2: a further rotation by the five high bits of $q$ \\
A permutation with no key, reading 30 of its 32 bits &
Anyone who knows the state can compute the rearrangement &
SBC version 2: a subkey for each permutation, and all 32 bits read \\
\hline
\end{tabular}
\caption{Weaknesses found in the October 2026 review and their repairs.}
\label{tab:fixes}
\end{table}

\subsection{Identical rounds and the slide attack}
\label{sec:slide}

In version 0 of SBC, step $i$ used the subkey $s_i = q(k_{i \bmod 16})$, where
$k_0, \ldots, k_{15}$ are the 16 words of the key. The subkeys therefore
repeat every 16 steps, the word positions repeat every 32, and \code{perm1},
which uses no key, runs every 64. Between one permutation and the next, the
cipher applied exactly the same function. A test confirmed that the state after
each permutation is that same function applied to the state after the one
before (Appendix \ref{sec:measure}).

A cipher built of identical keyed rounds is open to the \emph{slide attack} of
Biryukov and Wagner \cite{biryukov1999,biryukov2000}. Let $F$ be one round
under a fixed key, and let the cipher be $E = F^{r}$, that is $F$ applied $r$
times. Then
\begin{eqnarray}
E(F(P)) &=& F^{r+1}(P) \label{eq:slide_A} \\
        &=& F(E(P)). \label{eq:slide_B}
\end{eqnarray}
Suppose an attacker holds two plaintexts $P$ and $P'$ with $P' = F(P)$, a
\emph{slid pair}. By (\ref{eq:slide_B}) their ciphertexts satisfy
$C' = F(C)$. The attacker then has two inputs and outputs of a single round,
$(P, P')$ and $(C, C')$. A single round is far weaker than the whole cipher, so
the key can be sought from one round instead of 24.

The attacker cannot compute $F$, and so cannot make a slid pair on purpose, but
must find one by chance among many known plaintexts and ciphertexts. By the
birthday argument this takes about $2^{n/2}$ pairs for an $n$-bit block, and the
number does not depend on $r$. Adding rounds does not raise it.

For SBC, with $n = 1024$, the cost is about $2^{512}$ known blocks under one
key, which no one can collect; \code{fpwdman}, besides, uses a new key at every
save. The weakness is therefore one of principle, not of practice. It is still
a design defect, because it bypasses the number of rounds, which is the
cipher's main margin of safety.

\paragraph{The repair.}
Version 1 makes a separate subkey for every step from the key, by HKDF:
\begin{eqnarray}
(s_0, s_1, \ldots, s_{RK-1}) &=& \mathrm{HKDF}\bigl(\mathit{key},\
  \mathit{salt} = \text{empty}, \label{eq:v1_A} \\
 & & \quad \mathit{info} = \text{``SBC v1 subkeys''} \,\|\, K \,\|\, R,\
  L = 4RK \text{ bytes}\bigr), \nonumber
\end{eqnarray}
where $\|$ joins byte strings and $K$ and $R$ are written as four-byte numbers.
Each run of four output bytes is one subkey. Nothing else in the cipher
changes. No two rounds now use the same subkeys, and there is no simple rule
that links the subkeys of one key to those of another. Writing $K$ and $R$ into
the label keeps the schedules of different shapes unrelated: without them, the
24-round schedule would be the start of the 25-round one.

\paragraph{A repair that does not work.}
An obvious alternative is to add a different constant at each step, for
example
\begin{eqnarray}
s_i &=& q\bigl(k_{i \bmod 16} + (i+1)\varphi\bigr) \bmod 2^{32},
\label{eq:additive_A} \\
\varphi &=& \mathrm{9E3779B9}_{16}. \label{eq:additive_B}
\end{eqnarray}
Every step then has its own subkey, but the subkeys are related in a simple
way. Take a second key $k'$ whose every word is larger by $64\varphi$. Step $i$
under $k'$ then uses exactly the subkey of step $i+64$ under $k$. Sixty-four
steps is a whole number of rounds and also matches the permutation schedule, so
the two keys give the same cipher shifted by 64 steps, and the slide attack
returns for pairs of related keys. Any constant that grows regularly with the
step number has this flaw. HKDF has no such regularity. A milder alternative
that keeps the key word in every step is
\begin{equation}
s_i = q\bigl(k_{i \bmod N} \xor c_i\bigr),
\label{eq:table}
\end{equation}
with $c_i$ the first 32 bits of SHA-256 of ``SBC-RC'' and $i$. It also defeats
the slide attack, but it leaves the subkeys simply related to the key, so
version 1 uses HKDF.

\subsection{The unchecked iteration count}

The header of an FPMQ1 file states the number of PBKDF2 repetitions, and the
program used that number as written, checking only that it was positive. The
number cannot be authenticated before it is used, because the authentication
key itself comes out of PBKDF2. A crafted file giving 2,147,483,647 repetitions
would have made the program compute about 6.4 billion HMACs, roughly 40 minutes
in which the window would not respond. Nothing secret would leak, so the defect
was a nuisance rather than a break. Both formats now accept only 600,000, and
the check is made before any key is derived.

\subsection{The threefold cost of key derivation}

PBKDF2 produces 32 bytes for each full run of its repetitions. FPMQ1 asked for
96 bytes, so every open ran PBKDF2 three times: 1.8 million HMACs. But the
authentication key was the third 32-byte block by itself, so an attacker
testing a guess needed to compute only that block, 600,000 HMACs. The extra 1.2
million gave the user delay and no protection. FPMQ2 runs PBKDF2 once for 32
bytes and expands them with HKDF, which costs two HMACs. Saving and opening a
file now takes about a third as long as before (Appendix \ref{sec:measure}),
with the attacker's cost unchanged.

\subsection{The unauthenticated header}

FPMQ1 computed its authentication code over the encrypted text only. Changing
the salt or the repetition count would still have failed, because those change
the key, but the version name and the other header text were not protected.
FPMQ2 computes the code over both header lines and the encrypted text, and it
rejects any header that differs from what the program writes. In particular, an
FPMQ2 file relabelled as FPMQ1 is refused, so it cannot be steered onto the
older reading path. FPMQ3 does the same, and since each format derives its keys
with its own labels, an FPMQ3 file relabelled as FPMQ2 or FPMQ1 is refused too.

\subsection{The backup copy}

Before version 4.0, every save first copied the existing file to a backup with
the extension \code{.bak}. The backup kept the previous file exactly, in its old
format and under its old passphrase. Two cases made that a weakness. After a
legacy file was converted, the backup was the legacy file itself, which allows
fast guessing of a passphrase that also opens the new file. After a change of
master passphrase, the backup could still be opened with the old one, which
defeats the purpose of changing it. Version 4.0 keeps no backup. A failed save
is still harmless: the program writes the new file under a temporary name and
replaces the old one only when the new one is complete, so an interrupted save
leaves the old file as it was.

\subsection{The fixed rotation and the unkeyed permutation}
\label{sec:v2}

These two weaknesses were not found by an attack but by reasoning about where an
attack would begin; Section \ref{sec:open} (items \ref{item:perm} and
\ref{item:rotation}) gives the reasoning. Version 2 of SBC repairs both, and
FPMQ3 is the file format that uses it.

\emph{The rotation.} In versions 0 and 1, each output bit of $q$ depends only on
the input bits at its own place and below, and only the fixed rotation by seven
places carries information downward, over several steps. The step also keeps
the regular algebra of a quadratic. Version 2 adds the rotation of steps
(\ref{eq:stepv2_A}) to (\ref{eq:stepv2_C}): the changed word is rotated by the
five highest bits of $q$. A differential or linear attack must now predict one
of 32 rotation amounts at every step, and a rotation by a varying amount is not
a polynomial, so the step no longer has a quadratic's regular algebra. The
rotation is computed without any branch on its amount, so its timing reveals
nothing.

\emph{The permutation.} In versions 0 and 1, \code{perm1} uses no key: anyone who
knows the state at that point can compute the rearrangement. It also reads only
30 of the 32 bits of the number that drives it. Version 2 combines that number
with a secret subkey for each permutation, drawn from the same HKDF expansion as
the step subkeys, and rotates the number by five places between moves, so that
all 32 bits are read. Decryption is unchanged in form: it computes the number
from the rearranged block, combines it with the same subkey, and undoes the
moves in reverse order.

Version 2 needs 11 more subkeys than version 1, one for each permutation, and a
few more instructions per step. Its cost is negligible beside that of PBKDF2.

\subsection{Earlier repairs}

For completeness: the legacy format of the FLTK program had no key-derivation
function. The passphrase was scrambled into the key by a reversible function,
with no salt and no delay, and a wrong guess could be recognized after
decrypting a single block. Guessing was therefore millions of times cheaper
than for FPMQ1, FPMQ2 or FPMQ3. The Qt program corrected this with FPMQ1 before the
present review, and it now only reads the legacy format.

\section{Pseudo-code}
\label{sec:pseudo}

Throughout, words are 32 bits; $+$ and $-$ are modulo $2^{32}$; $\xor$ is
exclusive-or; $\rotl$ is left rotation. The cipher's parameters are the block
size $K$ (32), the key size $N$ (16 words), the rounds $R$ (24), and the
permutation interval $M = 2K$ (64).

\subsection{The basic SBC}

Algorithms \ref{alg:prim} to \ref{alg:cbc} give the cipher: its primitive
functions, its three key schedules, the encryption and decryption of one block,
and the chaining of blocks. The version $v$ (0, 1 or 2) selects the lines that
differ. Three points need comment. First, in \textsc{Perm1} the value $\sigma$
does not depend on the order of the words, so decryption can compute it from the
permuted block, and so undo the permutation. Second, in version 2 the rotation
amount $r$ depends only on $x_{p-1}$ and $s_i$, which the step leaves unchanged,
so decryption can compute it too. Third, the chaining starts from zero rather
than from a random value, but the first 16 bytes of every plaintext are random
(Algorithm \ref{alg:save}), which serves the same purpose.

\begin{algorithm}[htbp]
\caption{The primitive functions}
\label{alg:prim}
\begin{algorithmic}[1]
\Function{q}{$x$}
  \State \Return $(x+1)(2x+1)$ \Comment{modulo $2^{32}$}
\EndFunction
\Function{fold}{$j$} \Comment{stretch and fold of the indices $0 \ldots K-1$}
  \If{$2j \le K-1$}
    \State \Return $2j$
  \Else
    \State \Return $2(K-j)-1$
  \EndIf
\EndFunction
\Function{Perm1}{$x$, $\mathit{forward}$, $\kappa$, $\rho$}
  \Comment{versions 0, 1: $\kappa = 0$, $\rho = 7$; version 2: $\kappa = \pi_m$, $\rho = 5$}
  \State $\sigma \gets \kappa \xor q(x_0) \xor q(x_1) \xor \cdots \xor q(x_{K-1})$
    \Comment{does not depend on the order}
  \For{$t \gets 0$ \textbf{to} $7$}
    \State $\sigma_t \gets \sigma \rotl \rho t$
  \EndFor
  \If{$\mathit{forward}$}
    \For{$t \gets 0$ \textbf{to} $7$}
      \State for every $j$: $y_{\textsc{fold}((j + \sigma_t) \bmod K)} \gets x_j$;
             then $x \gets y$
    \EndFor
  \Else \Comment{undo the eight moves in reverse order}
    \For{$t \gets 7$ \textbf{down to} $0$}
      \State for every $j$: $y_j \gets x_{\textsc{fold}((j + \sigma_t) \bmod K)}$;
             then $x \gets y$
    \EndFor
  \EndIf
  \State \Return $x$
\EndFunction
\end{algorithmic}
\end{algorithm}

\begin{algorithm}[htbp]
\caption{Key setup}
\label{alg:keysetup}
\begin{algorithmic}[1]
\Function{KeySetupV0}{$k_0, \ldots, k_{N-1}$} \Comment{legacy and FPMQ1 only}
  \For{$i \gets 0$ \textbf{to} $RK-1$}
    \State $s_i \gets q(k_{i \bmod N})$ \Comment{repeats every $N$ steps}
  \EndFor
  \State \Return $s$
\EndFunction
\Function{KeySetupV1}{$\mathit{key}$} \Comment{any non-empty byte string}
  \State $\mathit{info} \gets$ ``SBC v1 subkeys'' $\|$ \textsc{be32}($K$) $\|$ \textsc{be32}($R$)
  \State $b \gets \textsc{HKDF-SHA-256}(\mathit{key}, \text{empty salt}, \mathit{info}, 4RK)$
  \For{$i \gets 0$ \textbf{to} $RK-1$}
    \State $s_i \gets$ bytes $4i \ldots 4i+3$ of $b$, read most significant first
  \EndFor
  \State \Return $s$ \Comment{$RK \le 2040$, the limit of HKDF}
\EndFunction
\Function{KeySetupV2}{$\mathit{key}$} \Comment{FPMQ3; any non-empty byte string}
  \State $P \gets \lfloor (RK-1)/M \rfloor$ \Comment{permutations per block: 11}
  \State $\mathit{info} \gets$ ``SBC v2 subkeys'' $\|$ \textsc{be32}($K$) $\|$ \textsc{be32}($R$)
  \State $b \gets \textsc{HKDF-SHA-256}(\mathit{key}, \text{empty salt}, \mathit{info}, 4(RK+P))$
  \State $w_0, \ldots, w_{RK+P-1} \gets$ the words of $b$, read most significant first
  \State $s_i \gets w_i$ for $0 \le i < RK$; \quad $\pi_m \gets w_{RK+m}$ for $0 \le m < P$
  \State \Return $s$, $\pi$ \Comment{$RK + P \le 2040$}
\EndFunction
\end{algorithmic}
\end{algorithm}

\begin{algorithm}[htbp]
\caption{Encrypting and decrypting one block}
\label{alg:block}
\begin{algorithmic}[1]
\Function{EncipherBlock}{$x_0, \ldots, x_{K-1}$; $s$, $\pi$, $v$}
  \For{$i \gets 0$ \textbf{to} $RK-1$}
    \State $p \gets i \bmod K$
    \State $t \gets q(x_{(p-1) \bmod K} \xor s_i)$
    \State $x_p \gets x_p + (t \rotl 7)$
    \If{$v = 2$}
      \State $x_p \gets x_p \rotl \lfloor t/2^{27} \rfloor$ \Comment{the five high bits of $t$}
    \EndIf
    \If{$i > 0$ \textbf{and} $i \bmod M = 0$}
      \If{$v = 2$}
        \State $x \gets \textsc{Perm1}(x, \text{true}, \pi_{i/M - 1}, 5)$
      \Else
        \State $x \gets \textsc{Perm1}(x, \text{true}, 0, 7)$
      \EndIf
    \EndIf
  \EndFor
  \State \Return $x$
\EndFunction
\Function{DecipherBlock}{$x_0, \ldots, x_{K-1}$; $s$, $\pi$, $v$}
  \For{$i \gets RK-1$ \textbf{down to} $0$}
    \If{$i > 0$ \textbf{and} $i \bmod M = 0$}
      \If{$v = 2$}
        \State $x \gets \textsc{Perm1}(x, \text{false}, \pi_{i/M - 1}, 5)$
      \Else
        \State $x \gets \textsc{Perm1}(x, \text{false}, 0, 7)$
      \EndIf
    \EndIf
    \State $p \gets i \bmod K$
    \State $t \gets q(x_{(p-1) \bmod K} \xor s_i)$
    \If{$v = 2$}
      \State $x_p \gets x_p \ggg \lfloor t/2^{27} \rfloor$ \Comment{right rotation}
    \EndIf
    \State $x_p \gets x_p - (t \rotl 7)$
  \EndFor
  \State \Return $x$
\EndFunction
\end{algorithmic}
\end{algorithm}

\begin{algorithm}[htbp]
\caption{Chaining blocks (CBC with a zero starting value)}
\label{alg:cbc}
\begin{algorithmic}[1]
\Function{EncipherCBC}{$P_0, \ldots, P_{n-1}$; $S$}
  \Comment{$S$: the schedule ($s$, $\pi$, $v$)}
  \State $C_{-1} \gets 0$
  \For{$b \gets 0$ \textbf{to} $n-1$}
    \State $C_b \gets \textsc{EncipherBlock}(P_b \xor C_{b-1}; S)$
  \EndFor
  \State \Return $C_0, \ldots, C_{n-1}$
\EndFunction
\Function{DecipherCBC}{$C_0, \ldots, C_{n-1}$; $S$}
  \State $C_{-1} \gets 0$
  \For{$b \gets 0$ \textbf{to} $n-1$}
    \State $P_b \gets \textsc{DecipherBlock}(C_b; S) \xor C_{b-1}$
  \EndFor
  \State \Return $P_0, \ldots, P_{n-1}$
\EndFunction
\end{algorithmic}
\end{algorithm}

\subsection{FPWDMAN}

Algorithm \ref{alg:save} saves the password file and Algorithm \ref{alg:open}
opens it. A legacy file is recognized by its first line, which holds two
numbers, and is decoded by a separate reader that reproduces the FLTK program
exactly. Saving any file, whatever its format, calls \textsc{Save}.

\begin{algorithm}[htbp]
\caption{Saving the password file (FPMQ3)}
\label{alg:save}
\begin{algorithmic}[1]
\Function{Save}{$\mathit{passphrase}$, $\mathit{entries}$}
  \State $\mathit{xml} \gets$ the entries written as XML
  \State $\mathit{salt} \gets$ 16 random bytes
  \State $\mathit{master} \gets \textsc{PBKDF2-HMAC-SHA-256}(\mathit{passphrase}, \mathit{salt}, 600000, 32)$
  \State $k_c \gets \textsc{HKDF-Expand}(\mathit{master}, \text{``FPMQ3 sbc''}, 64)$
  \State $k_m \gets \textsc{HKDF-Expand}(\mathit{master}, \text{``FPMQ3 mac''}, 32)$
  \State $F \gets$ 16 random bytes $\|$ \textsc{le32}(length of $\mathit{xml}$) $\|$ $\mathit{xml}$
  \State pad $F$ with bytes $\mathrm{FF}_{16}$ to a multiple of 128 bytes
  \State $C \gets \textsc{EncipherCBC}(F;\ \textsc{KeySetupV2}(k_c), v = 2)$
  \State $H \gets$ ``FPMQ3\textbackslash n'' $\|$
         ``kdf=pbkdf2-sha256; iters=600000; salt='' $\|$ \textsc{base64}($\mathit{salt}$) $\|$ ``\textbackslash n''
  \State $\mathit{tag} \gets \textsc{HMAC-SHA-256}(k_m, H \,\|\, C)$
  \State write $H$, then ``mac=hmac-sha256; tag='' $\|$ \textsc{base64}($\mathit{tag}$),
         then \textsc{base64}($C$) in lines of 64 characters
  \State erase $\mathit{master}$, $k_c$, $k_m$, $F$ and $\mathit{xml}$ from memory
\EndFunction
\end{algorithmic}
\end{algorithm}

\begin{algorithm}[htbp]
\caption{Opening the password file (FPMQ1, FPMQ2 or FPMQ3)}
\label{alg:open}
\begin{algorithmic}[1]
\Function{Open}{$\mathit{file}$, $\mathit{passphrase}$}
  \State read the version line, the salt, the repetition count $c$, the tag and $C$
  \If{the version is not FPMQ1, FPMQ2 or FPMQ3}
    \State \textbf{fail}: corrupt file
  \EndIf
  \If{$c \ne 600000$} \Comment{checked before any key is made}
    \State \textbf{fail}: corrupt file
  \EndIf
  \If{not FPMQ1 \textbf{and} the header differs from what \textsc{Save} writes}
    \State \textbf{fail}: corrupt file
  \EndIf
  \If{FPMQ3}
    \State derive $k_c$, $k_m$ as in \textsc{Save}; $\mathit{data} \gets H \,\|\, C$;
           $\mathit{schedule} \gets \textsc{KeySetupV2}$, $v = 2$
  \ElsIf{FPMQ2}
    \State the same, with the labels ``FPMQ2 sbc'' and ``FPMQ2 mac'';
           $\mathit{schedule} \gets \textsc{KeySetupV1}$, $v = 1$
  \Else \Comment{FPMQ1}
    \State $d \gets \textsc{PBKDF2-HMAC-SHA-256}(\mathit{passphrase}, \mathit{salt}, 600000, 96)$
    \State $k_c \gets d_{0 \ldots 63}$; $k_m \gets d_{64 \ldots 95}$; $\mathit{data} \gets C$;
           $\mathit{schedule} \gets \textsc{KeySetupV0}$, $v = 0$
  \EndIf
  \If{$\textsc{HMAC-SHA-256}(k_m, \mathit{data}) \ne \mathit{tag}$}
    \Comment{compared in constant time}
    \State \textbf{fail}: wrong passphrase, or the file was changed
  \EndIf
  \State $F \gets \textsc{DecipherCBC}(C;\ \mathit{schedule}(k_c))$
  \State skip 16 bytes; read the length; \Return the entries parsed from the XML
\EndFunction
\end{algorithmic}
\end{algorithm}

\FloatBarrier
\section{Remaining attacks and analyses}
\label{sec:open}

The review established the weaknesses above and measured some properties
(Appendix \ref{sec:measure}). It did not attempt the analyses below, which a
cipher must survive before it can be trusted for important data.

\subsection{Analyses of the cipher}

\begin{enumerate}
\item \textbf{Strength of a single round group.} The slide attack pays only
      if 64 steps are much weaker than 768. That is usual, but it has not been
      shown for SBC. How much work recovers the key from one input and output
      of 64 steps?

      A single step, by contrast, is easily undone, because $q$ is easy to
      invert. It is a T-function \cite{klimov2002}, and its derivative
      $q'(x) = 2nx + (an + c)$ is odd for every $x$. Once input bits 0 to
      $j-1$ are known, output bit $j$ is input bit $j$ plus a known quantity,
      so the input can be recovered from $y = q(x)$ bit by bit, from the
      lowest up, in 32 short steps. Newton's method is faster. It works modulo
      $2^{32}$ because $q'(x)$ is odd, and so has an inverse modulo $2^{32}$:
      \begin{eqnarray}
      x &\leftarrow& x - \bigl( q(x) - y \bigr)\, q'(x)^{-1} \bmod 2^{32}.
      \label{eq:qinv_A}
      \end{eqnarray}
      Each pass doubles the number of correct low bits, provided the starting
      value already has the right lowest bit. For $q(x) = (x+1)(2x+1)$ that
      bit is easy to find, because $q(x) \equiv x + 1 \pmod 2$: start from $y$
      with its lowest bit reversed. One correct bit then becomes 2, 4, 8, 16
      and 32, so five passes give all 32. (Tested on 20,000 random values:
      five passes inverted every one; four inverted about half.) Rivest characterized exactly which polynomials are one-to-one
      modulo $2^w$ \cite{rivest2001}, and inverting them is routine. Neither
      RC6 nor SBC relies on its quadratic being hard to invert: decryption
      never inverts $q$, but only computes it again, forward.

      The consequence is that one step whose input and output are both known
      gives up its subkey at once. If $x_p$ and $x_{p-1}$ are the words before
      a version 0 or 1 step and $x_p'$ the changed word after it, then
      \begin{eqnarray}
      t &=& (x_p' - x_p) \ggg 7, \label{eq:qinv_B} \\
      s_i &=& q^{-1}(t) \xor x_{p-1}. \label{eq:qinv_C}
      \end{eqnarray}
      In version 2 the rotation amount $r$ is unknown. An attacker can try
      all 32 values, compute $t = \bigl( (x_p' \ggg r) - x_p \bigr) \ggg 7$
      for each, and keep those for which $r$ equals the five high bits of $t$.
      The true value always survives, and each of the 31 wrong ones survives
      by chance with probability $1/32$, so on average about two candidates
      for the subkey remain. (Over 20,000 random steps the mean was 1.96, and
      the true subkey was always among them.) So in every version the
      protection lies not
      in any single step but in keeping the words between steps hidden.
      Within 64 steps each word is changed twice, and the question above is
      how much of that hiding survives.

\item \textbf{Differential cryptanalysis} \cite{biham1991}. Search for
      differences in the input that pass through a few rounds with high
      probability, and from them estimate how many rounds an attack could
      reach. The quadratic $q$ is a \emph{T-function} \cite{klimov2002}: each
      output bit depends only on the input bits at the same place and below.
      Differences in the low bits therefore propagate predictably within $q$,
      and only the rotation and the addition carry them upward.

\item \textbf{Linear cryptanalysis} \cite{matsui1993}. The two lowest output
      bits of $q$ are linear: bit 0 is the complement of input bit 0, and bit 1
      is a linear function of input bits 0 and 1. In version 0 of the 64-bit
      library a similar property made four bits of every step exactly linear.
      In the 32-bit cipher the addition in step (\ref{eq:step_B}) hides them
      behind carries. Whether useful linear approximations survive several
      rounds is not known.

\item \textbf{Reduced-round attacks.} Attack versions of SBC with 1, 2, 3, and
      more rounds, to learn how large the safety margin of 24 rounds actually
      is. This is the standard way of measuring a new cipher.

\item \label{item:perm}\textbf{The permutation.} In versions 0 and 1,
      \code{perm1} uses no key, and its choice depends on the exclusive-or of
      $q$ applied to every word, which does not depend on their order. It reads
      only 30 of those 32 bits. It produces a great variety of rearrangements
      (65,535 different ones from 65,536 random blocks), so it mixes well, but
      anyone who knows the state at that point can compute it. Version 2
      (Section \ref{sec:v2}) combines the driving number with a secret subkey
      for each permutation and reads all 32 bits. Questions remain for both
      forms. Does the permutation add strength, or only spreading? Are there
      blocks it leaves in a predictable order? The rearrangements all come from
      one family, eight moves of adding and folding, so even when keyed there
      are at most $2^{32}$ of them for each permutation, out of the $32!$, or
      about $2^{117}$, possible orders. Whether that family has structure an
      attacker can use is not known. A permutation drawn from all $32!$ orders
      with the key would remove the question, at the cost of more code.

\item \textbf{Algebraic attacks.} Each step is a polynomial of low degree over
      the integers modulo $2^{32}$, mixed with exclusive-or. Methods that
      solve systems of such equations should be tried on reduced rounds.

\item \label{item:rotation}\textbf{A rotation that depends on the data.} In
      versions 0 and 1 every step rotates by the same fixed amount, seven
      places. RC5 \cite{rivest1995} and RC6 \cite{RRSY98} instead rotate by
      an amount taken from the data; in RC6 it is the five high bits of a
      quadratic of another word, and what is rotated is the word being
      changed. Version 2 of SBC does the same (Section \ref{sec:v2}): it
      rotates the changed word $x_p$ by the five high bits of
      $q(x_{p-1} \xor s_i)$, after the fixed rotation and the addition. An
      earlier idea, rotating the output of $q$ by its own high bits, was set
      aside: there the amount is tied to the very value it moves, whereas the
      changed word holds independent data. The step remains reversible,
      because $x_{p-1}$ and $s_i$ are unchanged by it. The change helps in
      three ways.

      First, it breaks the property that each output bit of $q$ depends
      only on the input bits at the same place and below. The high bits of
      $q$ depend on every input bit, and they decide where every bit of $x_p$
      lands, so within a single step the low bits of the result depend
      on the high bits of the input. In versions 0 and 1 that happens only
      through the fixed rotation, over several steps.

      Second, it breaks the regular algebra of the step. A rotation by a
      varying amount is not a polynomial over the integers modulo $2^{32}$,
      so the regularities of a quadratic, such as its constant second
      difference, $q(x+d_1+d_2) - q(x+d_1) - q(x+d_2) + q(x) = 2n\,d_1 d_2$,
      no longer pass through the step intact. That makes the algebraic and
      higher-order differential attacks above harder.

      Third, a differential or linear attack follows a chosen pattern of
      differences or bits through the steps. With a data-dependent rotation
      it must also predict the rotation amount at every step, one of 32
      values, unless it can force those five bits. That multiplies its cost
      at every step.

      Data-dependent rotations have weaknesses of their own: the best attacks
      on RC5 work by controlling the rotation amounts \cite{biryukov1998}. What
      remains is to learn whether an attacker can steer the five high bits of
      $q$ often enough to matter, and to repeat for version 2 the differential,
      linear and reduced-round analyses of items 2 to 4.

\item \textbf{Statistical tests.} Run the output of SBC in counter mode
      through standard test batteries for random numbers, such as NIST SP
      800-22 or PractRand, at full and at reduced rounds.

\item \textbf{Related keys under versions 1 and 2.} With subkeys from HKDF, no
      relation between keys is known to carry over to the subkeys. In
      \code{fpwdman} the keys come from PBKDF2 and cannot be chosen by an
      attacker, so this matters mainly for SBC as a standalone cipher.
\end{enumerate}

\subsection{Analyses of the file format and the program}

\begin{enumerate}[resume]
\item \textbf{Old copies of the file.} A legacy copy of the file, under the same
      passphrase as the current one, lets an attacker guess that passphrase
      quickly and then open the current file. Version 4.0 makes no backup copy,
      but copies can survive elsewhere: a \code{.bak} written by an earlier
      version, the user's own backups, synchronized folders, and snapshots kept
      by the operating system. After converting, the user should delete such
      copies where possible and then change the master passphrase, which makes
      any copy that cannot be found harmless to the current file.

\item \label{item:argon}\textbf{Memory-hard key derivation.} PBKDF2 can be computed cheaply in
      parallel on graphics processors. Argon2id \cite{rfc9106}, which also
      requires much memory, is now preferred \cite{owasp}. A future format
      could use it.

\item \textbf{Unbounded inputs.} The program does not yet limit the size of a
      file it will read, nor the size a legacy file may claim for its
      decompressed contents. A crafted file could therefore demand a very large
      allocation. Limits should be set before any decoding.

\item \textbf{The length of the file.} Encryption hides the contents but not
      their approximate size, which reveals roughly how many entries the file
      holds. Padding to a fixed size would hide it.

\item \textbf{Timing.} The rotations, including version 2's data-dependent
      one, are single processor instructions whose time does not depend on the
      amount. But \code{perm1} writes words to places that depend on the data,
      and in version 2 on a subkey as well. On a shared computer, the timing of
      memory caches could in principle reveal something about those places,
      and so about the subkey. This matters only to an attacker running
      programs on the same machine at the same time.

\item \textbf{Secrets left in memory.} The program erases keys and plaintext
      after use, but Qt's strings are copied freely, so some copies cannot be
      reached and erased. The clipboard, the operating system's swap file and
      crash dumps can also hold secrets.

\item \textbf{Word size.} SBC could use 64-bit words, with a 64-bit $q$, in
      place of 32-bit ones. This would not make the password file noticeably
      safer. The practical attack on a stolen file is guessing the passphrase,
      which costs the same whatever the word size, and the key and block
      sizes, which would stay at 512 and 1024 bits, are already far beyond
      exhaustive search. Inside the cipher the change would cut both ways.
      Only the two lowest bits of $q$ are linear, so the weak share of each
      word would fall from 2 bits in 32 to 2 in 64, and the longer carries of
      64-bit addition and multiplication would mix the high bits more. But $q$
      moves information only upward within a word, and in versions 0 and 1 only
      the fixed rotation brings it down again: with a rotation of seven places,
      a high bit takes about nine steps to reach the bottom of a 64-bit word,
      against about five in a 32-bit word. (Version 2's data-dependent rotation
      narrows that difference.) A 1024-bit block would also hold 16 words instead of
      32, and the regularities of the quadratic, such as its constant second
      difference, would be unchanged. A 64-bit SBC would therefore have to be
      measured afresh, and it would be a new version of the cipher and of the
      file format.

\item \label{item:standard}\textbf{Comparison with standard ciphers.} SBC has had one reviewer. The
      standard ciphers, such as AES-256-GCM \cite{fips197} and
      XChaCha20-Poly1305, have had decades of study by many. Two earlier
      reviews of \code{fpwdman} recommended writing new files with one of them
      and reading SBC files only. That remains the most conservative choice;
      the work in this article is the alternative of strengthening SBC and
      inviting its analysis.
\end{enumerate}

\paragraph{Where further effort would count most.} Version 2 has made two of
the changes this list once recommended: the rotation that depends on the data
(item \ref{item:rotation}) and the keyed permutation (item \ref{item:perm}). In
order of their benefit to the safety of the password file, the most useful
remaining steps are these.
\begin{enumerate}[label=(\alph*)]
\item Argon2id in place of PBKDF2 (item \ref{item:argon}). It defends against
      the attack that matters most in practice, guessing the passphrase on
      many processors at once.
\item Analysis of version 2: the differential, linear and reduced-round
      studies of items 2 to 4, repeated for the new step, and the open
      questions of items \ref{item:perm} and \ref{item:rotation}.
\item A permutation drawn from all $32!$ orders with the key (item
      \ref{item:perm}), if that analysis finds the add-and-fold family to be a
      weakness.
\item A standard cipher for new files (item \ref{item:standard}). It gives
      the greatest assurance, at the price of giving up SBC.
\end{enumerate}
The first three keep SBC; the fourth replaces it. A change of word size is not
on the list, for the reasons given above. For now the author keeps SBC, so as to
continue its study.

\appendix

\section{Measurements}
\label{sec:measure}

All measurements below used keys and data generated at random for the purpose,
or fixed test vectors. A Python implementation of SBC version 0 was first
checked bit for bit against the compiled C++ program on a two-block test vector;
the version 1 and version 2 test vectors in the program's tests were computed
with a Python implementation written from the specification, using Python's own
HMAC and SHA-256, which also reproduces the version 0 vector.

\paragraph{Spreading of changes.} Mean number of the 1024 output bits changed
when one bit of the input is changed, over 100 random keys and blocks (a random
result would change 512):

\begin{center}
\begin{tabular}{rrr}
\hline
\textbf{Rounds} & \textbf{Mean bits changed} & \textbf{Fewest in any trial} \\
\hline
1  & 249.2 & 14 \\
2  & 512.2 & 481 \\
3  & 511.9 & 476 \\
4  & 513.4 & 470 \\
24 & 513.6 & 475 \\
\hline
\end{tabular}
\end{center}
Changing one bit of the key changed 510.5 bits on average after 24 rounds.

The same comparison for versions 1 and 2, over 2,000 random keys and blocks
each, with the standard error of the mean:

\begin{center}
\begin{tabular}{rrr}
\hline
\textbf{Rounds} & \textbf{Version 1} & \textbf{Version 2} \\
\hline
1 & $261.3 \pm 3.3$ & $266.5 \pm 3.3$ \\
2 & $511.8 \pm 0.4$ & $511.0 \pm 0.4$ \\
\hline
\end{tabular}
\end{center}
After one round the result depends strongly on which word was changed, since a
change late in the ring passes through few steps; in some trials only 2 or 3
bits changed. Across the block, then, version 2 spreads a change no faster than
version 1, which already matched a random result after two rounds. What version
2 changes is the behaviour within each word, below.

\paragraph{Within one step (version 2's rotation).} For random $x_p$, $x_{p-1}$
and $s_i$, one input bit of $x_{p-1}$ was flipped and the bits of the new $x_p$
that changed were counted, over 5,000 trials for each bit:

\begin{center}
\begin{tabular}{lrrrrr}
\hline
\textbf{Input bit flipped} & \textbf{0} & \textbf{8} & \textbf{16} & \textbf{24} & \textbf{31} \\
\hline
Version 1, bits changed & 16.45 & 13.42 & 9.19 & 5.18 & 2.02 \\
Version 2, bits changed & 15.95 & 15.92 & 15.76 & 15.57 & 16.06 \\
\hline
\end{tabular}
\end{center}
In version 1 a high input bit reaches few output bits, because each output bit
of $q$ depends only on the input bits at its own place and below. In version 2
every input bit changes about half of the 32. When bit 31 was flipped, over
20,000 trials, each of the output bits 0 to 5 changed with a frequency between
0.495 and 0.505 in version 2, and never in version 1.

\paragraph{Linearity of $q$.} The algebraic degree of output bits 0 to 5 of $q$,
as functions of the input bits at the same place and below, is 1, 1, 2, 3, 4,
5. Only the two lowest bits are linear.

\paragraph{The permutation.} In versions 0 and 1, \code{perm1} reads 30 of the
32 bits of its combined value (bits 9 and 10 are never used); with the rotation
by five of version 2 it reads all 32. Among 65,536 random values it produced
65,535 different rearrangements in one sample, and 65,536 in a second sample
for each of the two rotations.

In version 2, one random block rearranged under two random permutation subkeys
put on average 31.0 of its 32 words in different places, and in 2,000 trials the
two orders were never the same. Two blocks that differed in a single bit,
rearranged under the same subkey, never received the same order in 2,000 trials.

\paragraph{Timing of the rotation.} In the optimized program, the version 2
rotation compiles to two instructions, a shift that extracts the amount
(\code{shr ecx,1Bh}) and a rotation by it (\code{rol edx,cl}), with no branch.

\paragraph{Identical rounds in version 0.} In all 10 checks, the state after
permutation $m+1$ equalled one fixed 64-step function applied to the state after
permutation $m$. Under version 1, no two of the 24 rounds have the same
subkeys, while under version 0 all 24 do; the program's tests check both.

\paragraph{Speed.} One save and one open, in the program's debug build: 6.6
seconds for FPMQ1 and 2.3 seconds for FPMQ2. In a later run, made while other
programs were being compiled, FPMQ1 took 9.0 seconds, FPMQ2 3.0 and FPMQ3 3.0;
version 2 of the cipher adds no measurable time.

\paragraph{Tests.} Version 4.1 passes 41 tests of the cipher and key
derivation, including the published test vectors for SHA-256, HMAC-SHA-256,
PBKDF2 and HKDF and known-answer vectors for all three cipher versions, and 26
tests of the file formats, including the conversion of legacy, FPMQ1 and FPMQ2
files to FPMQ3, the rejection of altered FPMQ3 files, and the rejection of FPMQ3
files relabelled as older versions.

\clearpage
\bibliographystyle{ieeetr}
\bibliography{../../../CPVI/shared-tex/fdp}

\end{document}

% ---------------------------------------------
% Copyright Ben Paul Wise. All rights reserved.
% ---------------------------------------------
